\section{Conclusion}
\label{sec:conclusion}
In this work, we presented a scaling analysis of rectified flow models for
text-to-image synthesis. We proposed a novel timestep sampling for rectified
flow training that improves over previous diffusion training formulations for
latent diffusion models and retains the favourable properties of rectified
flows in the few-step sampling regime. We also demonstrated the advantages of
our transformer-based \modelname architecture that takes the multi-modal nature of the
text-to-image task into account. Finally, we performed a scaling study of this
combination up to a model size of 8B parameters and $5\times 10^{22}$ training
FLOPs. We showed that validation loss improvements correlate with both existing
text-to-image benchmarks as well as human preference evaluations. This, in
combination with our improvements in generative modeling and scalable, multimodal
architectures achieves performance that is competitive with state-of-the-art
proprietary models. 
The scaling trend shows no signs of saturation, which makes us optimistic
that we can continue to improve the performance of our models in the future.
%
%
